% META: {"id":"1SPE-PRODSCAL-CO-040","chapitre":"1SPE-PRODUIT-SCALAIRE","type_objet":"corrige","exercice_id":"1SPE-PRODSCAL-EX-040","status":"approved","origin":"generated","status_history":["generated","audited","accepted","approved"],"release_acceptance":"RELEASE_OWNER_BATCH_ACCEPTANCE","acceptance_closure_digest":"sha256:9b3ccf9a81c5520fbb7e03b7b2d3e7bf2057a02834b6d908bf3be5f49f3b553f","acceptance_receipt_digest":"sha256:4fad86f0282e0fa4e0419cca1ad6379ae7af5a280c04a4a90880f90a1c044242"}
% BEGIN-VERIFY
% from sympy import *
% # Isoceles triangle A(0,0), B(6,0), C(3,4)
% # Mediatrice de AB: x = 3
% # C is on mediatrice: x_C = 3 ✓, so CA = CB (isoceles)
% assert sqrt(9+16) == 5
% assert sqrt(9+16) == 5
% END-VERIFY
\begin{corrige}{1SPE-PRODSCAL-EX-040}
\begin{enumerate}
\item $CA = \sqrt{9 + 16} = 5$ et $CB = \sqrt{9 + 16} = 5$. Donc $CA = CB$ : triangle isocele en $C$.
\item La mediatrice de $[AB]$ a pour équation $x = 3$ (milieu $(3;0)$, $\overrightarrow{AB}$ horizontal). Or $x_C = 3$, donc $C$ est sur la mediatrice.
\item $\mathcal{A} = \dfrac{1}{2} \times AB \times h = \dfrac{1}{2} \times 6 \times 4 = 12$.
\item La hauteur issue de $C$ est perpendiculaire a $(AB)$ (axe $Ox$), donc $H = (3; 0)$ (le pied est le projete de $C$ sur $(AB)$).
\end{enumerate}
\end{corrige}
